\section*{Práce s ukázkou}

\subsection*{1. Zasazení ukázky do kontextu díla}
Ukázka pochází ze začátku knihy, představuje krysaře a Agnes a uvádí čtenáře so do děje.

\subsection*{2. Časoprostor}
\begin{itemize}
    \item \textbf{Čas:} Není napřímo zmíněný pravděpodobně středověk.
	\item \textbf{Prostor:} město Hammeln a jeho okolí
\end{itemize}

\subsection*{3. Téma a motivy}

\subsection*{4. Charakteristika postav}

\subsection*{5. Kompoziční výstavba}

\subsection*{6. Vypravěč a forma}

\subsection*{7. Typy promluv}

\subsection*{8. Veršová výstavba}

\subsection*{9. Jazyková rovina}

\subsection*{10. Umělecké prostředky, tropy a figury}

\subsection*{11. Vlastní názor na dílo}

\subsection*{12. Filmová a divadelní zpracování}

\newpage
